\section{模型范式变迁：从 GPU 到 Agent}

本章覆盖第一条主线：模型范式。故事从 1999 年第一颗 GPU 讲起，不是因为 GPU 是论文，而是因为它改变了可计算边界。Brook、CUDA、AlexNet、ResNet、Transformer 等工作，都说明一个朴素事实：算法范式能否成为时代主线，常常取决于它有没有抱住硬件和数据的大腿。

\begin{figure}[H]
\centering
\includegraphics[width=0.96\textwidth]{model-paradigm-timeline.png}
\caption{模型范式时间线：GPU、CNN、Attention、Transformer、MoE 与 Agent 逐步接力。自制概念图，依据 00:19:35--01:52:58 对谈内容整理。}
\end{figure}

\begin{knowledgebox}{读图：范式不是突然出现的}
AlexNet 之前有 GPU 和图像数据，Transformer 之前有 seq2seq、Attention、残差网络和并行计算，Agent 之前有 CoT、工具调用和 ReAct。范式变化常常是许多旧积累在某个硬件/数据窗口里突然连起来。
\end{knowledgebox}

\subsection{GPU、Brook、AlexNet 与硬件彩票}

本节先看硬件如何改变模型命运。Brook for GPU 代表早期通用 GPU 计算探索，AlexNet 则在 2012 年用 GPU 和深度卷积网络击穿 ImageNet。讲者反复强调“硬件彩票”：一个方法如果正好适合主流硬件，训练速度、工程生态和研究投入都会被放大；反之，思想可能正确，但难以成为时代主线。

\begin{figure}[H]
\centering
\includegraphics[width=0.94\textwidth]{hardware-lottery.png}
\caption{硬件彩票：能抱住主流硬件的架构会获得长期复利。自制概念图，依据 00:19:35--00:30:00 与 03:56:38--04:03:00 对谈内容整理。}
\end{figure}

\begin{knowledgebox}{读图：为什么 AlexNet 是分界线}
AlexNet 不只是“用了 CNN”，还证明了深度学习、GPU、数据集和工程技巧可以组合成压倒性优势。它击败的是手工特征主导的范式，也预示了 Bitter Lesson：当算力继续增长，通用学习方法会反复压过人工设计。
\end{knowledgebox}

\subsection{seq2seq、Attention 与 Transformer}

前面讲图像，本节转向序列。seq2seq 让模型把输入序列压缩成状态再解码，Attention 则允许解码时直接关注输入的不同位置，缓解长句和上下文丢失问题。Transformer 更进一步，用 self-attention 和并行计算替代 RNN 的顺序瓶颈，让序列内部 token 之间可以直接建模关系。

\begin{figure}[H]
\centering
\includegraphics[width=0.96\textwidth]{transformer-origin-ladder.png}
\caption{从序列建模到 Transformer：seq2seq、Attention、残差与并行计算共同铺路。自制概念图，依据 00:34:02--01:05:24 对谈内容整理。}
\end{figure}

\begin{importantbox}{Attention 的核心直觉}
Attention 允许模型在处理当前位置时直接计算它和其他位置的相关关系。相比 RNN 依赖一步步传递隐藏状态，Attention 缩短了信息路径；相比卷积需要多层扩大感受野，self-attention 可以在一层内看到整个上下文窗口。
\end{importantbox}

\subsection{ResNet、蒸馏、AlphaGo Zero 与 MoE}

前面讲的是架构如何处理图像和序列，本节把范围扩展到训练技巧、能力迁移和计算经济学。模型范式不仅是结构，还包括训练和能力迁移。ResNet 用残差连接解决更深网络训练难题；蒸馏提出大模型知识可以被小模型学习；AlphaGo Zero 展示强化学习和自我对弈可以绕开人类棋谱；现代 MoE 则通过专家路由降低激活成本，让大模型容量和计算成本不再完全线性绑定。

\begin{knowledgebox}{术语消化：模型范式节点}
\begin{tabularx}{\textwidth}{p{0.20\textwidth}p{0.32\textwidth}X}
\toprule
节点 & 解决的问题 & 后续影响 \\
\midrule
ResNet & 深层网络退化和训练困难 & 残差成为深度网络通用构件，也进入 Transformer。 \\
Distillation & 大模型能力如何迁移到小模型 & 支撑端侧部署、学生模型和模型压缩。 \\
AlphaGo Zero & 不依赖人类数据也能学习策略 & 强化学习、自我博弈和 test-time search 的思想被反复引用。 \\
MoE & 增加参数容量但控制每次计算量 & DeepSeek 等模型把专家路由变成现代 LLM 的重要路线。 \\
\bottomrule
\end{tabularx}
\end{knowledgebox}

\subsection{CoT、LoRA、ReAct 与 Agent 化}

本节从基础模型进入使用方式。CoT 让大家意识到模型输入会影响输出，推动 prompt engineering 和 context engineering；LoRA 让低成本参数适配变得普遍；ReAct 把 reasoning 和 acting 放在同一个循环中，为 Agent 从理论走向工具调用奠基。

\begin{figure}[H]
\centering
\includegraphics[width=0.96\textwidth]{adaptation-agent-stack.png}
\caption{从能力到行动的栈：蒸馏、LoRA、CoT 和 ReAct 把模型推向可用系统。自制概念图，依据 01:26:40--01:37:10 对谈内容整理。}
\end{figure}

\begin{knowledgebox}{读图：模型能力需要被“引出”和“接住”}
蒸馏和 LoRA 解决能力迁移与适配成本，CoT 让模型显式展开推理，ReAct 让推理连接外部工具。Agent 并不是凭空出现，而是这些机制逐渐把模型推向行动闭环。
\end{knowledgebox}

\subsection{Bitter Lesson：长期主义的冷水}

前面所有节点都在暗示一个规律，本节用 The Bitter Lesson 把它说透。它的观点是，长期看，依赖通用计算和大规模搜索/学习的方法，往往胜过手工编码知识的方法。讲者强调，这不是说手工特征永远无用，而是说手工结构常常在某个算力量级内有效；一旦算力和数据跨过阈值，通用学习方法会重新占优。

\begin{figure}[H]
\centering
\includegraphics[width=0.94\textwidth]{bitter-lesson-map.png}
\caption{Bitter Lesson：长期看，通用计算与规模化学习会反复压过手工结构。自制概念图，依据 01:45:00--01:52:40 对谈内容整理。}
\end{figure}

\begin{warningbox}{Bitter Lesson 不是“不要做结构设计”}
结构设计仍然重要，尤其在数据少、算力弱、目标明确时。但如果一个结构不能随算力、数据和硬件生态一起 scale，它就可能在下一轮范式迁移中被通用方法压过。
\end{warningbox}

\subsection{本章小结}

模型范式线告诉我们：AI 进步不是单篇论文胜利，而是硬件、数据、训练方法、架构和使用方式共同对齐。GPU 让深度学习可行，Attention 让长依赖可学，Transformer 抱住并行硬件，CoT/ReAct 则把模型推向行动。

